\documentclass[a4paper,10pt]{article}

\usepackage[utf8]{inputenc}
\usepackage[T1]{fontenc}
\usepackage[french]{babel}
\usepackage{geometry}
\geometry{a4paper, margin=1cm}

\usepackage{hyperref}
\hypersetup{
    colorlinks=true,
    linkcolor=blue,
    urlcolor=blue,
}

\usepackage{enumitem}
\setlist{nosep, left=0pt}

\usepackage{titlesec}
\titleformat{\section}{\large\bfseries}{}{0em}{}[\titlerule]
\titlespacing{\section}{0pt}{10pt}{5pt}

\usepackage{parskip}
\setlength{\parskip}{0.5em}

\begin{document}

% ===== EN-TÊTE =====
\begin{center}
    {\Huge \textbf{Nathan LEFEBVRE}} \\[5pt]
    \textbf{Développement Logiciel} \\[10pt]
    \begin{tabular}{c}
        \href{mailto:nathan.lefebvre@utbm.fr}{nathan.lefebvre@utbm.fr} \\
        06 43 03 00 91 \\
        \href{https://linkedin.com/in/nathan-lefebvre-2b1a11327}{linkedin.com/in/nathan-lefebvre} \\
        22/06/2006 \\[5pt]
        Belfort, Bourgogne-Franche-Comté
    \end{tabular}
\end{center}

% ===== PROFIL =====
\section{Profil}
Actuellement en 2\textsuperscript{e} année de cycle préparatoire à l'UTBM, je réalise un semestre d'étude en Chine qui renforce mon autonomie et mon ouverture d'esprit. Curieux et rigoureux, j'ai un intérêt marqué pour la compréhension des systèmes et la résolution de problèmes logiques. Je recherche aujourd'hui une alternance en développement logiciel afin de mettre en pratique mes compétences et de progresser sur des projets concrets en informatique.

% ===== FORMATION =====
\section{Formation}
\textbf{École d'ingénieur UTBM} – Sevenans (90) \\
2\textsuperscript{e} année cycle préparatoire · Semestre d'étude en Chine

\textbf{Baccalauréat (mention bien)} – Lycée Cassini, Clermont (60) \\
Spécialités NSI, Mathématiques

% ===== EXPÉRIENCES PROFESSIONNELLES =====
\section{Expériences professionnelles}

\subsection*{Stage ouvrier – EJ Picardie (fonderie)}
Saint-Crépin-lbouvilliers (60)
\begin{itemize}
    \item Gestion de production : utilisation d'un logiciel de planification pour organiser les tâches, suivre les séries de fabrication et affecter la production aux opérateurs.
    \item Fabrication de noyaux : production sur machines industrielles.
    \item Ligne de production : insertion manuelle des noyaux dans les demi-moules.
    \item Polyvalence multi-ateliers : intervention sur différents postes selon les besoins.
    \item Suivi opérationnel : participation aux réunions quotidiennes.
\end{itemize}

\subsection*{Travail saisonnier – UCAC \& VALFRANCE}
Avrigny, Cambronne-lès-Clermont, Charny (77)
\begin{itemize}
    \item Analyse de céréales : mesure de l'humidité, du poids spécifique.
    \item Saisie et structuration des données dans un logiciel interne.
    \item Suivi des flux logistiques : réception et expédition de camions.
    \item Organisation du stockage des céréales selon leurs caractéristiques.
    \item Travail en environnement soutenu.
    \item Polyvalence terrain : entretien et utilisation d'engins.
\end{itemize}

\subsection*{Stage de découverte – Valfrance}
Senlis (60)
\begin{itemize}
    \item Observation du déploiement d'un logiciel de suivi de maintenance.
    \item Découverte de l'organisation interne et des différents pôles.
    \item Participation à des réunions de suivi.
    \item Réalisation d'une cartographie de l'organisation.
    \item Première immersion en environnement professionnel.
\end{itemize}

% ===== COMPÉTENCES TECHNIQUES =====
\section{Compétences techniques}
\textbf{C} (Intermédiaire – SDL) \quad
\textbf{Python} (Intermédiaire) \quad
\textbf{SQL} (notions) \quad
\textbf{HTML/CSS} (notions) \quad
\textbf{PHP} (notions) \quad
\textbf{Linux – Ubuntu} (notions)

% ===== LANGUES =====
\section{Langues}
Français (natif) \quad
Anglais (B1) \quad
Mandarin (notions) \quad
Espagnol (notions)

% ===== CENTRES D'INTÉRÊT =====
\section{Centres d'intérêt}
Développement de jeux vidéo (projets personnels) \quad
Musculation, course à pied \quad
Lecture : fantaisie, romans, mangas

\end{document}